\newpage
\section{Integrator}\label{sec:1}

\subsection{Entwurf der Integratorschaltung}\label{sec:1-1}
\subsubsection{Aufgabenstellung}
Mittels einer Integratorschaltung soll für eine konstante Eingangspannung $U_e=\SI{100}{\milli\volt}$ eine Ausgangsspannungsänderung von $\SI{500}{\milli\volt}$ erzeugt werden.

\subsubsection{Versuchsaufbau und Dimensionierung}

\begin{figure}[h]
	\begin{center}
		\includegraphics[width=0.55\textwidth]{Integrator_Schaltung.png}
	\end{center}
	\caption{Umkehrintegrator}\label{fig:integrator-schaltung}
\end{figure}

Wenn man $I_N=0$ annimmt, kann man sich daraus $\tau$ berechnen. Die Ausgangsspannungsänderung $\frac{d U_A}{dt}$ muss negativ sein, da $RC$ positiv ist.

\begin{gather}
	I_1=-I_C \\
	\frac{U_1}{R}=-C \frac{d U_A}{dt} \\
	RC = - \frac{U_1}{\frac{d U_A}{dt}} = - \frac{\SI{100}{\milli\volt}}{\SI{-500}{\milli\volt\per\second}} = \SI{0.2}{\second} \\
	\tau = RC = \SI{0.2}{\second} \label{eq:1-tau}
\end{gather}

Wird für den Kondensator ($C$) $\SI{1}{\mu\farad}$ gewählt so ergibt sich aus Gl.~\ref{eq:1-tau} für  den Widerstand ($R$) $\SI{200}{\kilo\ohm}$. Der Widerstand wurde mit der R-Dekade RS-200 siehe Geräteliste~\ref{sec:geraeteliste} realisiert.


